\appendix

\section{Instance Generation Parameters}
\label{app:instances}

Table~\ref{tab:instance_params} documents the full parameterization of our synthetic instance generator.
All instances are generated with master seed 42 for reproducibility.

\begin{table}[h]
\centering
\caption{Synthetic instance generation parameters.}
\label{tab:instance_params}
\small
\begin{tabular}{lccccccc}
\toprule
\textbf{Config} & $|\mathcal{S}|$ & $|\mathcal{A}|$ & $|\mathcal{C}|$ & $|\mathcal{I}|$ & $|\mathcal{L}|$ & $T$ \\
\midrule
S10  & 10  & 5   & 1  & 2  & 3  & 15  \\
S15  & 15  & 8   & 1  & 2  & 4  & 25  \\
S20  & 20  & 10  & 2  & 3  & 5  & 35  \\
S25  & 25  & 12  & 2  & 4  & 6  & 45  \\
S30  & 30  & 15  & 2  & 4  & 7  & 55  \\
S50  & 50  & 20  & 3  & 6  & 10 & 90  \\
S75  & 75  & 35  & 4  & 8  & 12 & 120 \\
S100 & 100 & 43  & 5  & 10 & 15 & 150 \\
\bottomrule
\end{tabular}
\end{table}

For each instance, parameters are drawn as follows:

\paragraph{Scene parameters.}
\begin{itemize}
    \item Scene durations: $\delta_s \sim U\{1, 2, 3, 4\}$ hours;
    \item Each scene requires 1--4 actors, assigned with uniform probability;
    \item Each scene is compatible with 1--2 locations;
    \item Time windows assigned to $\sim$40\% of scenes;
    \item Soft deadlines assigned to $\sim$30\% of scenes, penalties $\rho_s \sim U(50, 500)$.
\end{itemize}

\paragraph{Actor parameters.}
\begin{itemize}
    \item Daily wages: $w_a \sim U(80, 1000)$ (following~\citet{liu2019});
    \item Availability: $\sim$80\% of the planning horizon;
    \item Criticality weights: $\kappa^A_a \sim U(1, 10)$;
    \item Maximum consecutive working days: $\bar{h}_a \sim U\{3, 4, 5, 6\}$.
\end{itemize}

\paragraph{Location hierarchy parameters.}
\begin{itemize}
    \item Cities: $|\mathcal{C}| = \max(1, \lfloor |\mathcal{L}|/3 \rfloor)$;
    \item Towns: $|\mathcal{I}| = \max(2, \lfloor 2|\mathcal{L}|/3 \rfloor)$, distributed across cities;
    \item Location availability: $\sim$85\% of the planning horizon;
    \item Criticality weights: $\kappa^L_l \sim U(1, 10)$;
    \item Intra-town transfer costs: $\tau^{\text{loc}} \sim U(100, 500)$;
    \item Inter-town transfer costs: $\tau^{\text{town}} \sim U(500, 2000)$;
    \item Inter-city transfer costs: $\tau^{\text{city}} \sim U(2000, 8000)$.
\end{itemize}

\paragraph{Equipment parameters (C14).}
\begin{itemize}
    \item Equipment types: $|\mathcal{E}| = \max(2, \lfloor |\mathcal{L}|/2 \rfloor)$;
    \item Equipment supply: $\xi_e \sim U\{1, 2\}$ units per type;
    \item Scene equipment requirements: $\sim$50\% of scenes require 1--2 equipment types.
\end{itemize}

\paragraph{Actor booking parameters (C15).}
\begin{itemize}
    \item Booking windows: $\sim$60\% of actors have booking window $[\beta^s_a, \beta^e_a]$, window length $\sim U(T/4, 3T/4)$;
    \item Holding fees: $h_a \sim U(0.2, 0.6) \times w_a$ (20--60\% of daily wage).
\end{itemize}

\paragraph{Scene coordination parameters (C16--C17).}
\begin{itemize}
    \item Co-scheduling groups: $|\mathcal{G}| = \max(1, \lfloor n/8 \rfloor)$ groups of 2--3 scenes each;
    \item Separation pairs: $|\mathcal{F}| = \max(1, \lfloor n/10 \rfloor)$ pairs with gap $\delta_{ij} \sim U(2, T/10)$.
\end{itemize}

\paragraph{Location and lead-actor parameters (C18--C19).}
\begin{itemize}
    \item Location concurrency capacity: $\mu_l \sim U\{1, 2, 3\}$;
    \item Lead actors: $|\mathcal{A}^P| = \max(1, \lfloor m/5 \rfloor)$ (20\% of cast);
    \item Rest periods: $r_a \sim U\{1, 2\}$ days for $a \in \mathcal{A}^P$.
\end{itemize}

\paragraph{Structural parameters.}
\begin{itemize}
    \item Precedence constraints: $\sim$20\% of scene pairs form a sparse DAG;
    \item Daily capacity: $W_{\max} = 8$ hours;
    \item All 19 constraint families (C1--C19) are active.
\end{itemize}

\section{Hyperparameter Sensitivity Analysis}
\label{app:sensitivity}

We performed a sensitivity analysis of key BDM-AOS hyperparameters on the S25 instance.
For each parameter, we varied one parameter while fixing all others at their default values.

\paragraph{Population Size.}
Varying $N \in \{10, 20, 30, 50\}$: cost improves modestly from $N=10$ to $N=20$ (3.2\% average reduction) but shows diminishing returns beyond $N=20$, while runtime scales linearly with $N$.

\paragraph{Generations.}
Varying $G \in \{20, 40, 60, 100\}$: convergence is achieved by generation 30--35 for most instances, confirming that $G=40$ provides sufficient search budget.

\paragraph{Q-Learning Parameters.}
The AOS learning rate $\alpha \in \{0.05, 0.1, 0.2\}$ and discount factor $\gamma \in \{0.8, 0.9, 0.95\}$ show stable performance across tested ranges, indicating robustness of the Q-learning adaptation.

\section{Q-Learning State-Action Table}
\label{app:qtable}

The AOS state space discretizes four continuous features:
\begin{itemize}
    \item \textbf{Generation progress}: early ($g/G < 0.33$), mid ($0.33 \leq g/G < 0.67$), late ($g/G \geq 0.67$);
    \item \textbf{Population diversity}: low ($d < 0.3$), medium ($0.3 \leq d < 0.7$), high ($d \geq 0.7$);
    \item \textbf{Improvement rate}: stagnant ($r < 0.01$), slow ($0.01 \leq r < 0.1$), fast ($r \geq 0.1$);
    \item \textbf{Quality gap}: good ($q < 0.1$), moderate ($0.1 \leq q < 0.5$), poor ($q \geq 0.5$).
\end{itemize}

This yields $3^4 = 81$ states.
The action space comprises $3 \times 3 = 9$ operator pairs: $\{$PMX, OX, LAX$\} \times \{$swap, block-swap, ACM$\}$.
The Q-table is initialized to zeros and updated via standard Q-learning (Eq.~\ref{eq:qlearning}).

\section{Extended Results Tables}
\label{app:extended_results}

Full per-seed results for all 8 algorithms across all 8 instances are provided in the supplementary repository at \url{https://github.com/VJlaxmi/movie-scheduling-paper/tree/main/experiments/runs}.
Each JSON file contains the complete run record including final cost, makespan, runtime, Pareto front, convergence history, and AOS operator frequency log.

\section{Friedman Test Details}
\label{app:friedman}

We apply the Friedman test across all algorithms on each instance, using mean cost over 30 seeds as the performance measure.
When the Friedman test rejects the null hypothesis of equal ranks ($p < 0.05$), we apply Holm's step-down procedure for pairwise comparisons with BDM-AOS as the control algorithm.

\section{Decomposition Analysis}
\label{app:decomposition}

The SIG-based Louvain decomposition produces the following block structures:

\begin{table}[h]
\centering
\caption{Block decomposition statistics across instances.}
\label{tab:decomposition}
\small
\begin{tabular}{lccc}
\toprule
\textbf{Config} & \textbf{Num. Blocks} & \textbf{Avg. Block Size} & \textbf{Search Space Reduction} \\
\midrule
S10  & 3   & 3.3 & $10!/3! = 604{,}800\times$   \\
S15  & 4   & 3.8 & $15!/4! = 5.4 \times 10^{10}\times$  \\
S20  & 6   & 3.3 & $20!/6! = 3.4 \times 10^{15}\times$  \\
S25  & 6   & 4.2 & $25!/6! = 2.2 \times 10^{22}\times$  \\
S30  & 8   & 3.8 & $30!/8! = 8.1 \times 10^{28}\times$  \\
S50  & 12  & 4.2 & $50!/12! = 1.3 \times 10^{55}\times$ \\
S75  & 16  & 4.7 & $75!/16! = 5.4 \times 10^{96}\times$ \\
S100 & 24  & 4.2 & $100!/24! = 4.1 \times 10^{134}\times$ \\
\bottomrule
\end{tabular}
\end{table}

\section{Runtime Environment}
\label{app:runtime}

\begin{table}[h]
\centering
\caption{Software environment for all experiments.}
\label{tab:environment}
\small
\begin{tabular}{ll}
\toprule
\textbf{Component} & \textbf{Version} \\
\midrule
Python           & 3.14.2 \\
NumPy            & 2.3.0 \\
SciPy            & 1.16.0 \\
pandas           & 2.3.0 \\
NetworkX         & 3.5 \\
python-louvain   & 0.16 \\
Pyomo            & 6.9.2 \\
CBC Solver       & 2.10 \\
Matplotlib       & 3.10.3 \\
OS               & Windows \\
\bottomrule
\end{tabular}
\end{table}
